% =============================================================================
% PREÁMBULO LATEX EXACTO - COPIAR LITERALMENTE
% =============================================================================
\documentclass[12pt,a4paper]{report}

% Fuentes CRÍTICAS para compatibilidad Word
\usepackage[utf8]{inputenc}
\usepackage[T1]{fontenc}
\usepackage{helvet}           % Helvetica para texto principal
\usepackage{courier}           % Courier para código
\usepackage{graphicx}
\usepackage{fancyhdr}
\usepackage{hyperref}
\usepackage{listings}
\usepackage{xcolor}
\usepackage{booktabs}
\usepackage{geometry}
\usepackage{microtype}
\usepackage{tocbibind}
\usepackage{amsmath}
\usepackage{amssymb}
\usepackage{appendix}
\usepackage{caption}
\usepackage{subcaption}
\usepackage{enumitem}
\usepackage{titlesec}

% Márgenes
\geometry{
    left=2.5cm,
    right=2.5cm,
    top=2.5cm,
    bottom=2.5cm,
    headheight=15pt
}

% Listings para código
\lstset{
    basicstyle=\ttfamily\fontfamily{pcr}\selectfont,
    backgroundcolor=\color{gray!10},
    frame=single,
    frameround=tttt,
    numbers=left,
    numberstyle=\tiny\color{gray},
    keywordstyle=\color{blue}\bfseries,
    commentstyle=\color{green!60!black},
    stringstyle=\color{red},
    breaklines=true,
    tabsize=2,
    captionpos=b,
    xleftmargin=2em,
}

% Headers y footers
\pagestyle{fancy}
\fancyhf{}
\fancyhead[L]{\color{blue!70!black}\textbf{\sffamily\leftmark}}
\fancyhead[R]{\color{blue!70!black}\textbf{\sffamily\thepage}}
\fancyfoot[C]{\color{gray}\small\sffamily Project Name}

% Títulos con sans-serif
\titleformat{\chapter}[display]
    {\normalfont\huge\bfseries\sffamily\color{blue!70!black}}
    {\chaptertitlename\ \thechapter}{20pt}{\Huge}

\titleformat{\section}
    {\normalfont\Large\bfseries\sffamily\color{blue!70!black}}
    {\thesection}{1em}{}

\titleformat{\subsection}
    {\normalfont\large\bfseries\sffamily\color{orange!80!black}}
    {\thesubsection}{1em}{}

% Hyperref
\hypersetup{
    colorlinks=true,
    linkcolor=blue!70!black,
    urlcolor=blue!70!black,
    pdftitle={Project Name},
    pdfauthor={wisrovi},
}

% =============================================================================
% REGLAS DE ORO
% =============================================================================
% 1. USAR SOLO helvet (\sffamily) + courier (\ttfamily)
% 2. EVITAR lmodern, fuentes matemáticas complejas
% 3. CAPÍTULOS: \cleardoublepage al inicio
% 4. TÍTULOS: \sffamily en todos
% 5. CÓDIGO: \fontfamily{pcr}\selectfont
% 6. IMÁGENES: Solo PNG/JPG, rutas a sources/
